\vspace{2.5pt plus 3pt minus 0.5pt}\noindent\hspace*{1.2em}\kbvnum{(5)}\hspace{0.55em}“ನೀನು?” ಎಂಬ ಸರ್ವನಾಮವನ್ನು ಪುಲ್ಲಿಂಗ, ಸ್ತ್ರೀಲಿಂಗ, ಹಾಗೂ ನಪುಂಸಕ ಲಿಂಗಗಳಲ್ಲಿ ಬಳಸಬಹುದೆಂಬುದನ್ನು ಉದಾಹರಣೆಗಳ ಮೂಲಕ ಸಮರ್ಥಿಸಿ.\par

\vspace{2.5pt plus 3pt minus 0.5pt}\noindent\hspace*{1.2em}\kbvnum{(6)}\hspace{0.55em}ನಪುಂಸಕ ಲಿಂಗಕ್ಕೆ ಒಂದು ಉದಾಹರಣೆಯನ್ನು ಕೊಡಿ.\par

\vspace{2.5pt plus 3pt minus 0.5pt}\noindent\hspace*{1.2em}\kbvnum{(7)}\hspace{0.55em}ನಿತ್ಯ ನಪುಂಸಕ ಲಿಂಗ ಎಂದರೇನು ? ಎರಡು ಉದಾಹರಣೆಗಳನ್ನು ಕೊಡಿ.\par

\vspace{2.5pt plus 3pt minus 0.5pt}\noindent\hspace*{1.2em}\kbvnum{(8)}\hspace{0.55em}“ನಾನು ಈ ದೇಶದ ಪ್ರಜೆ.”— ಈ ವಾಕ್ಯದಲ್ಲಿರುವ ನಾಮಪದಗಳನ್ನು ಗುರ್ತಿಸಿ.ಲಿಂಗವನ್ನು ತಿಳಿಸಿ.\par

\kbvchapter{ಪಾಠ ೧೦}{ನಾಮಪದಗಳ ರೂಪ ಬದಲಾವಣೆ ({\kbvlatin ii})}

\kbvsubtitle{ವಚನ}

\indent ಒಂದು ಗಾಡಿ ಮುಂದೆ ಹೋಗುತ್ತಿದೆ ; ಆದರೆ ಹಿಂದೆ ಅನೇಕ ಗಾಡಿಗಳು ಬರುತ್ತಿವೆ.\par

\indent ಈ ವಾಕ್ಯಗಳಲ್ಲಿ ದಪ್ಪಕ್ಷರಗಳಲ್ಲಿರುವ “ಗಾಡಿ” ಗಾಡಿಗಳು” ಎಂಬ ಶಬ್ದಗಳನ್ನು ಗಮನಿಸಿ. ಇಲ್ಲಿ “ಗಾಡಿ” ಎಂಬ ಶಬ್ದ ಒಂದು ವಸ್ತುವನ್ನು ಮಾತ್ರ ಸೂಚಿಸುತ್ತದೆ ; ಆದರೆ “ಗಾಡಿಗಳು” ಎಂಬ ಶಬ್ದ ಒಂದಕ್ಕಿಂತ ಹೆಚ್ಚು ವಸ್ತುಗಳು ಎಂಬುದನ್ನು ಸೂಚಿಸುತ್ತದೆ.\par

\indent ಹೀಗೆ, ಒಂದು ವಸ್ತುವನ್ನು ಕುರಿತಾದ ಶಬ್ದವನ್ನು ಏಕವಚನ ಶಬ್ದ ಎಂದು ಒಂದಕ್ಕಿಂತ ಹೆಚ್ಚು ವಸ್ತುಗಳನ್ನು ಕುರಿತಾದ ಶಬ್ದವನ್ನು ಬಹುವಚನ ಶಬ್ದವೆಂದು ಹೇಳುತ್ತೇವೆ. ಕೆಳಗಿನ ಉದಾಹರಣೆಗಳನ್ನು ನೋಡಿ...\par

\vspace{5pt plus 4pt minus 1pt}
\begin{center}
{\renewcommand{\arraystretch}{1.24}\setlength{\tabcolsep}{10pt}\rowcolors{2}{kbvrow}{white}
\begin{tabularx}{0.88\textwidth}{X X}
{\kbvsecfont\color{kbvaccentdark}ಏಕವಚನ} & {\kbvsecfont\color{kbvaccentdark}ಬಹುವಚನ} \\[1pt]
\arrayrulecolor{kbvline}\specialrule{0.5pt}{0pt}{2pt}\arrayrulecolor{black}
ಆಳು & ಆಳುಗಳು \\
ಹಸು & ಹಸುಗಳು \\
ರೂಪಾಯಿ & ರೂಪಾಯಿಗಳು \\
ಗಿಣಿ & ಗಿಣಿಗಳು \\
\end{tabularx}}
\end{center}
\vspace{5pt plus 4pt minus 1pt}
